% ============================================================
% Proposta de Projeto — Disciplina de Projeto de Sistemas
% Ubíquos e Embarcados — UFSC Araranguá — 2026.1
% ============================================================
\documentclass[12pt,a4paper]{article}

% ---- Codificação e idioma ----
\usepackage[utf8]{inputenc}
\usepackage[T1]{fontenc}
\usepackage[brazil]{babel}

% ---- Margens e espaçamento (NBR-like) ----
\usepackage[top=3cm, bottom=2cm, left=3cm, right=2cm]{geometry}
\usepackage{setspace}
\onehalfspacing

% ---- Tipografia ----
\usepackage{lmodern}
\usepackage{microtype}

% ---- Links e referências ----
\usepackage[hidelinks]{hyperref}
\usepackage{url}

% ---- Figuras e tabelas ----
\usepackage{graphicx}
\usepackage{float}
\usepackage{booktabs}
\usepackage{array}
\usepackage{longtable}
\usepackage{caption}

% ---- TikZ para os diagramas ----
\usepackage{tikz}
\usetikzlibrary{shapes.geometric, arrows.meta, positioning, automata, calc, fit}

% ---- Listas ----
\usepackage{enumitem}

% ---- Cabeçalho / rodapé ----
\usepackage{fancyhdr}
\pagestyle{fancy}
\fancyhf{}
\rhead{\small Projeto de Sistemas Ubíquos e Embarcados — 2026.1}
\lhead{\small UFSC Araranguá}
\cfoot{\thepage}
\renewcommand{\headrulewidth}{0.4pt}

% ---- Seções sem numeração romana ----
\renewcommand{\thesection}{\arabic{section}}

% ============================================================
\begin{document}
\selectlanguage{brazil}

% ============================================================
% CAPA
% ============================================================
\begin{titlepage}
  \centering
  {\large Universidade Federal de Santa Catarina\\
         Centro de Ciências, Tecnologias e Saúde\\
         Campus Araranguá\\[2mm]
         Curso de Engenharia de Computação}

  \vspace{3cm}

  {\LARGE\bfseries Controle de Acesso Inteligente para\\
  Imóveis de Aluguel por Temporada\\
  Baseado em IoT com ESP32 e MQTT}

  \vspace{1.5cm}

  {\large\textbf{Modelo de Proposta de Projeto}}\\[4mm]
  {\large Projeto de Sistemas Ubíquos e Embarcados — 2026.1}

  \vfill

  \begin{flushleft}
    \textbf{Alunos:}\\[2mm]
    1.\ Fernando Doqui Futila
  \end{flushleft}

  \vspace{2cm}
  {\large Araranguá, maio de 2026}
\end{titlepage}

% ============================================================
% SUMÁRIO
% ============================================================
\tableofcontents
\newpage

% ============================================================
\section{Introdução}
% ============================================================

\subsection{Contextualização e Relevância do Tema}

O fenômeno da economia compartilhada, consolidado globalmente por plataformas como
Airbnb e Booking.com, transformou unidades residenciais convencionais em ativos de
hotelaria dinâmica. Essa transição, no entanto, impõe desafios de infraestrutura
significativos, uma vez que os sistemas de controle de acesso tradicionais não possuem
a granularidade necessária para gerenciar hóspedes de curta duração.

A automação residencial, antes vista apenas como um item de luxo, tornou-se uma
ferramenta de gestão essencial para garantir a viabilidade operacional e a segurança
desses imóveis. Neste contexto, a Internet das Coisas (IoT) surge como o elo
tecnológico capaz de resolver essa lacuna, permitindo a digitalização da barreira
física através de dispositivos conectados à rede mundial de computadores.

Os sistemas ubíquos e embarcados são protagonistas nesse cenário: microcontroladores
de baixo custo como o ESP32 combinam conectividade Wi-Fi, capacidade de processamento
local e interface com sensores físicos em um único chip, viabilizando soluções que
antes exigiam hardware proprietário de alto custo.

\subsection{Justificativa}

A justificativa deste trabalho fundamenta-se na carência de soluções que integrem,
em um único ecossistema de baixo custo, o controle de acesso temporário e o
monitoramento ambiental reativo. Embora existam protótipos acadêmicos e produtos
comerciais de alto custo, poucos conseguem atender simultaneamente aos requisitos de:
\begin{enumerate}
  \item Geração e expiração automática de credenciais vinculadas ao calendário de
        reservas;
  \item Gestão eficiente de múltiplos usuários e propriedades;
  \item Integração nativa com sensores de segurança (gás, fumaça e temperatura);
  \item Baixo custo de replicação através de hardware \emph{open-source}.
\end{enumerate}

A maioria das fechaduras inteligentes comerciais é voltada para uso residencial fixo,
apresentando limitações críticas para o modelo de aluguel: incapacidade de gerar
credenciais temporárias automáticas, necessidade de interação presencial para
configuração e ausência de integração com sensores ambientais que monitorem a
integridade do imóvel durante a estadia.

\subsection{Problemática}

O problema central reside na falta de interoperabilidade e na dependência de
conectividade constante. Sistemas baseados puramente em nuvem falham quando a rede
local do imóvel apresenta instabilidade, enquanto sistemas puramente locais (como
teclados de senha fixa) não permitem a revogação dinâmica de acessos.

Além disso, a segurança patrimonial em imóveis alugados é frequentemente negligenciada:
vazamentos de gás ou princípios de incêndio podem ocorrer sem que o proprietário seja
notificado ou que o sistema de acesso tome medidas automáticas para facilitar a
evacuação ou o socorro.

A gestão de chaves físicas ainda é uma realidade para muitos anfitriões, gerando:
\begin{itemize}
  \item Custos logísticos de deslocamento para entrega e recolhimento de chaves;
  \item Riscos de duplicação não autorizada;
  \item Total ausência de auditoria sobre quem acessa a propriedade e em quais horários.
\end{itemize}

\subsection{Objetivo Geral}

Este projeto tem como objetivo desenvolver e validar uma solução de controle de acesso
inteligente baseada em IoT, voltada especificamente para imóveis de aluguel por
temporada, integrando o gerenciamento remoto de chaves digitais e o monitoramento
preventivo de sensores ambientais, por meio de um dispositivo embarcado de baixo custo
(ESP32) conectado a uma plataforma \emph{backend} Django via protocolo MQTT sobre
HiveMQ Cloud.

% ============================================================
\section{Trabalhos Relacionados}
% ============================================================

A pesquisa sobre fechaduras inteligentes avançou significativamente com a
popularização de microcontroladores de baixo custo. Contudo, a aplicação voltada
especificamente para o mercado de locação por temporada ainda apresenta lacunas
operacionais. A seguir são apresentados os principais trabalhos relacionados
identificados na literatura.

\subsection*{Trabalho 1 — Smart Lock System Using BLE and Multi-Factor Authentication}

Proposta de arquitetura centrada em Bluetooth Low Energy (BLE) que exige autenticação
multifator para liberar a porta. Embora ofereça alta segurança contra clonagem de
sinal, a dependência da proximidade física do hóspede dificulta a gestão remota e a
auditoria em tempo real, uma vez que os registros não são sincronizados
instantaneamente com a nuvem. Não há integração com sensores ambientais.

\subsection*{Trabalho 2 — Temporary QR Code-Based Access Control for Short-Term Rentals}

Propõe o uso de QR Codes temporários gerados pelo servidor para abertura de portas.
A principal vantagem é a capacidade de operar \emph{offline}, mas a solução carece de
um mecanismo de revogação imediata: caso um hóspede precise ser impedido de entrar
antes do fim do prazo, o código físico ou digital em sua posse continuaria válido até
a expiração local. Não há monitoramento ambiental.

\subsection*{Trabalho 3 — Retrofit de Fechaduras Tradicionais com Microcontroladores e
Identificação por MAC Address}

Foca no \emph{retrofit} de fechaduras tradicionais, minimizando custos de instalação.
Entretanto, a segurança baseada apenas no endereço MAC do dispositivo móvel é
insuficiente para o mercado de aluguel profissional, dado que esse identificador pode
ser falsificado por softwares comuns, e não há suporte a múltiplos usuários com janelas
temporais distintas.

\subsection*{Trabalho 4 — IoT-Based Home Automation and Security System Using MQTT and ESP8266}

Sistema de automação residencial baseado em MQTT com ESP8266 que integra controle
de iluminação, temperatura e alarmes. Demonstra a viabilidade do protocolo MQTT
para ambientes com recursos limitados, mas não trata o ciclo de vida de credenciais
temporárias nem o conceito de "servidor interceptador" — os comandos podem ser
publicados diretamente no broker por qualquer cliente autenticado, sem validação de
regras de negócio.

\subsection*{Trabalho 5 — Secure Access Control System for IoT Using Token-Based Authentication
and Cloud Broker}

Propõe um sistema de controle de acesso IoT com autenticação por token JWT e broker
MQTT em nuvem. A solução é mais próxima da proposta deste projeto, porém não contempla
o monitoramento de sensores ambientais (gás, fumaça, temperatura) integrado ao mesmo
dispositivo, nem o fluxo de resgate de credencial pelo hóspede sem criação de conta
formal.

\bigskip
\noindent Ao comparar as soluções existentes, percebe-se que nenhuma integra
simultaneamente o controle de acesso temporal dinâmico com a segurança reativa baseada
em sensores ambientais. É nesta lacuna que o presente trabalho se posiciona, oferecendo
uma plataforma integrada que prioriza tanto a conveniência do hóspede quanto a
integridade do patrimônio do anfitrião.

% ============================================================
\section{Requisitos do Sistema}
% ============================================================

\subsection{Requisitos Funcionais}

\begin{longtable}{p{1.100cm} p{3.800cm} p{10.00cm}}
  \toprule
  \textbf{\#} & \textbf{Módulo} & \textbf{Descrição} \\
  \midrule
  \endfirsthead
  \toprule
  \textbf{\#} & \textbf{Módulo} & \textbf{Descrição} \\
  \midrule
  \endhead
  \bottomrule
  \endfoot

  RF01 & Autenticação & O sistema deve permitir que anfitriões se registrem com nome,
         e-mail e senha, e façam login com geração de token de autenticação. \\[3pt]

  RF02 & Gestão de fechaduras & O anfitrião deve poder cadastrar, editar e remover
         fechaduras eletrônicas, associadas por endereço MAC do ESP32. \\[3pt]

  RF03 & Abertura remota & O anfitrião deve poder acionar a abertura da fechadura
         remotamente pela interface web, mediante publicação de comando MQTT. \\[3pt]

  RF04 & Acesso temporário & O sistema deve gerar tokens de resgate de 12 caracteres
         hexadecimais, com janela temporal de validade definida (início e fim de reserva). \\[3pt]

  RF05 & Ciclo de vida do acesso & O status do acesso temporário deve seguir a máquina
         de estados: \texttt{PENDENTE} $\rightarrow$ \texttt{AGENDADO} $\rightarrow$
         \texttt{ATIVO} $\rightarrow$ \texttt{EXPIRADO} / \texttt{REVOGADO}. \\[3pt]

  RF06 & Resgate de token pelo hóspede & O hóspede deve poder resgatar o token via
         aplicativo mobile sem necessidade de criar conta no sistema. \\[3pt]

  RF07 & Abertura pelo hóspede & O sistema deve validar a janela temporal da reserva
         antes de publicar o comando de abertura, rejeitando acessos fora do período. \\[3pt]

  RF08 & Revogação de acesso & O anfitrião deve poder revogar um acesso temporário
         ativo a qualquer momento, impedindo novas aberturas imediatamente. \\[3pt]

  RF09 & Telemetria de sensores & O firmware deve publicar leituras dos sensores
         (temperatura, umidade, concentração de gás) a cada intervalo configurável. \\[3pt]

  RF10 & Monitoramento ambiental & O sistema deve registrar alertas automáticos quando
         o sensor de gás ultrapassar limiar de risco, atuando de forma autônoma no
         firmware (abertura preventiva do relé) sem depender de conectividade. \\[3pt]

  RF11 & Heartbeat e status online & O firmware deve publicar mensagem de heartbeat
         \texttt{"ONLINE"} a cada 30 segundos; a desconexão deve atualizar o campo
         \texttt{esta\_online} via Last Will Testament (LWT). \\[3pt]

  RF12 & Histórico de eventos & O sistema deve registrar todos os eventos de acesso,
         alertas de sensor e eventos de sistema em log auditável com timestamp. \\[3pt]

  RF13 & Interface web (dashboard) & Deve existir painel web responsivo para anfitriões
         com telas de fechaduras, reservas, logs e perfil. \\[3pt]

  RF14 & Aplicativo mobile & Deve existir aplicativo mobile para hóspedes (iOS e
         Android) para resgate de token e abertura da porta. \\[3pt]
\end{longtable}

\subsection{Requisitos Não Funcionais}

\begin{longtable}{p{1.100cm} p{3.800cm} p{10.00cm}}
  \toprule
  \textbf{\#} & \textbf{Categoria} & \textbf{Descrição} \\
  \midrule
  \endfirsthead
  \toprule
  \textbf{\#} & \textbf{Categoria} & \textbf{Descrição} \\
  \midrule
  \endhead
  \bottomrule
  \endfoot

  RNF01 & Segurança & Toda comunicação entre backend e broker MQTT deve utilizar TLS
          (porta 8883) com autenticação por usuário e senha. \\[3pt]

  RNF02 & Segurança & Nenhuma interface de usuário (web ou mobile) deve se comunicar
          diretamente com o broker MQTT ou com o hardware (padrão Servidor Interceptador). \\[3pt]

  RNF03 & Segurança & A API REST deve exigir token de autenticação em todos os
          endpoints de anfitrião; endpoints de hóspede autenticam-se pelo token\_resgate. \\[3pt]

  RNF04 & Precisão de dados & Valores de sensores devem ser preservados com no mínimo
          4 algarismos significativos em toda a cadeia: firmware $\rightarrow$ MQTT
          $\rightarrow$ banco de dados $\rightarrow$ interface. \\[3pt]

  RNF05 & Desempenho & A latência entre o acionamento do botão de abertura na interface
          e a atuação do relé no ESP32 deve ser inferior a 3 segundos em condições
          normais de rede. \\[3pt]

  RNF06 & Disponibilidade & O firmware deve operar de forma autônoma para eventos de
          emergência (detecção de gás) independentemente da disponibilidade da rede. \\[3pt]

  RNF07 & Escalabilidade & A arquitetura backend deve suportar múltiplos anfitriões,
          múltiplas fechaduras e múltiplos hóspedes simultâneos. \\[3pt]

  RNF08 & Portabilidade & O backend deve ser independente de motor de banco de dados,
          suportando SQLite3 (desenvolvimento) e PostgreSQL (produção) sem alterações
          no código-fonte. \\[3pt]

  RNF09 & Usabilidade & O dashboard web deve ser responsivo e funcional em resoluções
          a partir de 768px. O aplicativo mobile deve funcionar nos sistemas iOS e
          Android via Expo. \\[3pt]

  RNF10 & Rastreabilidade & Todo evento de acesso, alerta ou sistema deve ser
          persistido com timestamp em \texttt{HistoricoEvento} para fins de auditoria. \\[3pt]

  RNF11 & Custo de hardware & O custo total dos componentes de hardware por unidade
          instalada (ESP32 + sensores + relé) deve ser inferior a R\$100,00. \\[3pt]
\end{longtable}

% ============================================================
\section{Proposta / Desenvolvimento}
% ============================================================

\subsection{Descrição da Solução}

A proposta consiste no desenvolvimento de um ecossistema de controle de acesso
inteligente para imóveis de aluguel por temporada, composto por quatro camadas
integradas: firmware embarcado em ESP32, servidor backend Django, painel web
para anfitriões e aplicativo mobile para hóspedes.

O sistema adota a arquitetura denominada \textbf{Servidor Interceptador}: nenhuma
interface de usuário se comunica diretamente com o hardware. Todo comando de abertura
obrigatoriamente passa pelo servidor Django, que valida regras de negócio (identidade
do solicitante, janela temporal da reserva) antes de publicar o payload \texttt{"OPEN"}
no tópico MQTT correspondente ao dispositivo. O ESP32, que mantém subscrição nesse
tópico via HiveMQ Cloud (TLS 8883), recebe o comando e aciona o relé da fechadura.

O fluxo completo de um acesso temporário envolve:
\begin{enumerate}
  \item O anfitrião cadastra a fechadura com o endereço MAC do ESP32 pelo dashboard;
  \item O anfitrião cria uma reserva, definindo o e-mail do hóspede e o período;
  \item O sistema gera automaticamente um token de resgate de 12 caracteres hexadecimais;
  \item O hóspede insere o token no aplicativo, que valida junto ao servidor
        (status: \texttt{PENDENTE} $\to$ \texttt{AGENDADO});
  \item Dentro da janela reservada, o hóspede pressiona "Abrir Porta"; o servidor
        valida a janela temporal e publica \texttt{"OPEN"} no tópico MQTT;
  \item O ESP32 recebe o comando e aciona o relé; o evento é registrado no histórico.
\end{enumerate}

Paralelamente, o firmware publica telemetria de sensores (temperatura, umidade,
concentração de gás) a cada 30 segundos. Em caso de detecção de gás acima do limiar
de risco, o firmware atua de forma autônoma — abrindo o relé para evacuação — sem
depender da conectividade com o servidor.

\subsection{Diagrama Esquemático da Arquitetura do Sistema}

A Figura~\ref{fig:arq-sistema} apresenta o diagrama de arquitetura do sistema,
ilustrando as quatro camadas (interfaces, servidor, mensageria e hardware) e os
canais de comunicação entre elas.

\begin{figure}[H]
  \caption{Diagrama de arquitetura do sistema — padrão Servidor Interceptador.}
  \label{fig:arq-sistema}
  \centering
  \begin{tikzpicture}[>=Stealth,
    frontend/.style={rectangle, rounded corners=4pt, draw=blue!70!black,
      fill=blue!8, minimum width=3.8cm, minimum height=1.1cm,
      align=center, font=\small},
    backend/.style={rectangle, rounded corners=4pt, draw=teal!70!black,
      fill=teal!8, minimum width=6.5cm, minimum height=1.1cm,
      align=center, font=\small},
    broker/.style={rectangle, rounded corners=4pt, draw=orange!70!black,
      fill=orange!8, minimum width=6.5cm, minimum height=1.1cm,
      align=center, font=\small},
    hw/.style={rectangle, rounded corners=4pt, draw=red!70!black,
      fill=red!8, minimum width=6.5cm, minimum height=1.1cm,
      align=center, font=\small},
    sensor/.style={rectangle, rounded corners=2pt, draw=gray!60,
      fill=gray!6, minimum width=2.0cm, minimum height=0.85cm,
      align=center, font=\scriptsize},
    arr/.style={->, thick, draw=gray!70!black},
    bidarr/.style={<->, thick, draw=gray!70!black},
    lbl/.style={font=\scriptsize, fill=white, inner sep=1.5pt}
  ]
  % Nós
  \node[frontend] (app)    at (0.8,  6.0) {App do Hóspede\\{\scriptsize React Native / Expo}};
  \node[frontend] (dash)   at (6.0,  6.0) {Dashboard Web\\{\scriptsize React 19 + TypeScript}};
  \node[backend]  (django) at (3.4,  4.0) {Django REST API\\{\scriptsize Servidor Interceptador + paho-mqtt}};
  \node[broker]   (mqtt)   at (3.4,  2.0) {HiveMQ Cloud\\{\scriptsize Broker MQTT $\cdot$ TLS 8883}};
  \node[hw]       (esp32)  at (3.4,  0.0) {Firmware ESP32\\{\scriptsize WiFiClientSecure + PubSubClient}};
  \node[sensor]   (relay)  at (1.0, -1.9) {Relé\\{\tiny GPIO 26}};
  \node[sensor]   (mq2)    at (3.4, -1.9) {Sensor MQ-2\\{\tiny GPIO 34}};
  \node[sensor]   (dht)    at (5.8, -1.9) {Sensor DHT22\\{\tiny GPIO 27}};
  % Setas interface → backend
  \draw[arr] (app.south)  -- ++(0,-0.35) -| ([xshift=-1.0cm]django.north)
      node[pos=0.28, lbl, left] {\texttt{HTTPS/REST}};
  \draw[arr] (dash.south) -- ++(0,-0.35) -| ([xshift=+1.0cm]django.north)
      node[pos=0.28, lbl, right] {\texttt{HTTPS/REST}};
  % Setas backend ↔ broker ↔ ESP32
  \draw[bidarr] (django) -- (mqtt)
      node[midway, lbl, right=4pt] {\texttt{publish / subscribe}};
  \draw[bidarr] (mqtt)   -- (esp32)
      node[midway, lbl, right=4pt] {\texttt{commands / status / sensors}};
  % Setas ESP32 → sensores/atuadores
  \draw[arr] (esp32.south) -| (relay.north);
  \draw[arr] (esp32.south)  -- (mq2.north);
  \draw[arr] (esp32.south) -| (dht.north);
  % Linhas divisoras de camada
  \draw[dashed, gray!35] (-1.0, 5.05) -- (8.2, 5.05);
  \draw[dashed, gray!35] (-1.0, 3.05) -- (8.2, 3.05);
  \draw[dashed, gray!35] (-1.0, 1.05) -- (8.2, 1.05);
  \draw[dashed, gray!35] (-1.0,-0.95) -- (8.2,-0.95);
  % Rótulos de camada
  \node[font=\tiny\itshape, text=gray!70, anchor=east] at (-1.1, 5.5)  {Interfaces};
  \node[font=\tiny\itshape, text=gray!70, anchor=east] at (-1.1, 4.0)  {Servidor};
  \node[font=\tiny\itshape, text=gray!70, anchor=east] at (-1.1, 2.0)  {Mensageria};
  \node[font=\tiny\itshape, text=gray!70, anchor=east] at (-1.1, 0.0)  {Hardware};
  \end{tikzpicture}
  \par\smallskip{\footnotesize Fonte: Próprio autor.}
\end{figure}

\subsection{Diagrama Entidade-Relacionamento do Banco de Dados}

A Figura~\ref{fig:der} apresenta o Diagrama Entidade-Relacionamento das três entidades
centrais que sustentam a lógica do sistema: \texttt{Fechadura}, \texttt{AcessoTemporario}
e \texttt{HistoricoEvento}.

\begin{figure}[H]
  \caption{Diagrama Entidade-Relacionamento do banco de dados.}
  \label{fig:der}
  \centering
  \begin{tikzpicture}[>=Stealth,
    eheader/.style={rectangle, draw=black, fill=gray!15,
      minimum width=5.2cm, font=\small\bfseries, align=center,
      inner xsep=4pt, inner ysep=3pt},
    eattr/.style={rectangle, draw=gray!50, fill=white,
      minimum width=5.2cm, font=\footnotesize, text ragged,
      inner xsep=6pt, inner ysep=2pt},
    epk/.style={rectangle, draw=gray!50, fill=white,
      minimum width=5.2cm, font=\footnotesize\bfseries, text ragged,
      inner xsep=6pt, inner ysep=2pt},
    efk/.style={rectangle, draw=gray!50, fill=white,
      minimum width=5.2cm, font=\footnotesize\itshape, text ragged,
      inner xsep=6pt, inner ysep=2pt},
    rel/.style={thick, draw=black},
    card/.style={font=\small}
  ]
  %% Fechadura
  \node[eheader] (fh)  at (0.0,  0.0)  {Fechadura};
  \node[epk,  below=-\pgflinewidth of fh] (f1) {\underline{id} (PK)};
  \node[eattr,below=-\pgflinewidth of f1] (f2) {id\_unico (UUID)};
  \node[eattr,below=-\pgflinewidth of f2] (f3) {nome};
  \node[eattr,below=-\pgflinewidth of f3] (f4) {id\_dispositivo (MAC)};
  \node[efk,  below=-\pgflinewidth of f4] (f5) {proprietario (FK → User)};
  \node[eattr,below=-\pgflinewidth of f5] (f6) {esta\_online};
  \node[eattr,below=-\pgflinewidth of f6] (f7) {ultima\_comunicacao};
  %% AcessoTemporario
  \node[eheader] (ah)  at (7.0,  0.0)  {AcessoTemporario};
  \node[epk,  below=-\pgflinewidth of ah] (a1) {\underline{id} (PK)};
  \node[efk,  below=-\pgflinewidth of a1] (a2) {fechadura (FK → Fechadura)};
  \node[eattr,below=-\pgflinewidth of a2] (a3) {token\_resgate (12 hex)};
  \node[eattr,below=-\pgflinewidth of a3] (a4) {hospede\_identificador};
  \node[eattr,below=-\pgflinewidth of a4] (a5) {inicio\_reserva};
  \node[eattr,below=-\pgflinewidth of a5] (a6) {fim\_reserva};
  \node[eattr,below=-\pgflinewidth of a6] (a7) {status (5 estados)};
  %% HistoricoEvento
  \node[eheader] (hh)  at (0.0, -5.6)  {HistoricoEvento};
  \node[epk,  below=-\pgflinewidth of hh] (h1) {\underline{id} (PK)};
  \node[efk,  below=-\pgflinewidth of h1] (h2) {fechadura (FK → Fechadura)};
  \node[eattr,below=-\pgflinewidth of h2] (h3) {tipo (ACESSO, ALERTA...)};
  \node[eattr,below=-\pgflinewidth of h3] (h4) {valor\_sensor (Decimal, 4 dp)};
  \node[eattr,below=-\pgflinewidth of h4] (h5) {descricao};
  \node[eattr,below=-\pgflinewidth of h5] (h6) {timestamp};
  %% Relacionamentos
  \draw[rel] (fh.east) -- (ah.west)
      node[pos=0.05, above, card] {1}
      node[pos=0.95, above, card] {N};
  \draw[rel] (fh.south) -- (hh.north)
      node[pos=0.05, right, card] {1}
      node[pos=0.95, right, card] {N};
  \end{tikzpicture}
  \par\smallskip{\footnotesize Fonte: Próprio autor.}
\end{figure}

\subsection{Diagrama de Estados do Acesso Temporário}

A Figura~\ref{fig:estados} apresenta a máquina de estados do \texttt{AcessoTemporario},
evidenciando as cinco transições possíveis ao longo do ciclo de vida de uma reserva.

\begin{figure}[H]
  \caption{Máquina de estados do \texttt{AcessoTemporario}.}
  \label{fig:estados}
  \centering
  \begin{tikzpicture}[>=Stealth, shorten >=2pt,
    every state/.style={minimum size=1.5cm, font=\scriptsize\bfseries,
      align=center, draw=black!70},
    trlbl/.style={font=\tiny, fill=white, inner sep=1.5pt, align=center}
  ]
  \node[state, initial left, fill=yellow!20]  (pend) at (0,   0) {PENDENTE};
  \node[state, fill=blue!12]                  (agen) at (3.8, 0) {AGENDADO};
  \node[state, fill=green!15]                 (ativ) at (7.6, 0) {ATIVO};
  \node[state, accepting, fill=gray!15]       (expi) at (11.4,0) {EXPIRADO};
  \node[state, accepting, fill=red!12]        (revo) at (5.7,-3) {REVOGADO};
  % Fluxo principal
  \draw[->] (pend) -- (agen)
      node[midway, above, trlbl] {hóspede resgata\\token};
  \draw[->] (agen) -- (ativ)
      node[midway, above, trlbl] {primeiro acesso\\dentro da janela};
  \draw[->] (ativ) -- (expi)
      node[midway, above, trlbl] {prazo\\expirado};
  % Revogação
  \draw[->, bend right=25] (pend) to
      node[midway, left,  trlbl] {anfitrião\\revoga} (revo);
  \draw[->] (agen.south) -- (revo.north)
      node[midway, left,  trlbl] {anfitrião\\revoga};
  \draw[->, bend left=25] (ativ)  to
      node[midway, right, trlbl] {anfitrião\\revoga} (revo);
  \end{tikzpicture}
  \par\smallskip{\footnotesize Fonte: Próprio autor.}
\end{figure}

\subsection{Diagrama de Sequência — Abertura da Porta pelo Hóspede}

A Figura~\ref{fig:sequencia} ilustra o fluxo de mensagens durante uma operação de
abertura de porta solicitada pelo hóspede, desde o acionamento no aplicativo até
a atuação do relé no ESP32.

\begin{figure}[H]
  \caption{Diagrama de sequência da abertura da porta pelo hóspede.}
  \label{fig:sequencia}
  \centering
  \begin{tikzpicture}[>=Stealth, font=\small]
    % Linhas de vida
    \foreach \x/\lbl in {
        0/App Hóspede, 3/Django API, 6/HiveMQ Cloud, 9/ESP32} {
      \node[draw, minimum width=2.4cm, minimum height=0.7cm,
            rounded corners=2pt, align=center, font=\scriptsize\bfseries]
            at (\x, 0) {\lbl};
      \draw[dashed, gray!50] (\x,-0.35) -- (\x,-7.0);
    }
    % Mensagens
    \draw[->] (0,-0.9) -- (3,-0.9)
        node[midway, above, font=\scriptsize] {POST /hospede/abrir/};
    \draw[->] (3,-1.6) -- (3,-1.6)
        node[right, font=\scriptsize] {Valida token\_resgate e janela temporal};
    \draw[->] (3,-2.3) -- (6,-2.3)
        node[midway, above, font=\scriptsize] {publish(``OPEN'')};
    \draw[->] (6,-3.0) -- (9,-3.0)
        node[midway, above, font=\scriptsize] {deliver topic: v1/locks/\{MAC\}/commands};
    \draw[->] (9,-3.7) -- (9,-3.7)
        node[right, font=\scriptsize] {Aciona relé (GPIO 26) por 3 s};
    \draw[->] (9,-4.4) -- (6,-4.4)
        node[midway, above, font=\scriptsize] {publish status: ``ONLINE''};
    \draw[->] (6,-5.1) -- (3,-5.1)
        node[midway, above, font=\scriptsize] {deliver topic: v1/locks/\{MAC\}/status};
    \draw[->] (3,-5.8) -- (3,-5.8)
        node[right, font=\scriptsize] {Cria HistoricoEvento (ACESSO)};
    \draw[->] (3,-6.5) -- (0,-6.5)
        node[midway, above, font=\scriptsize] {HTTP 200 OK};
  \end{tikzpicture}
  \par\smallskip{\footnotesize Fonte: Próprio autor.}
\end{figure}

\subsection{Ferramentas de Hardware e Software}

\subsubsection{Hardware}

\begin{table}[H]
  \centering
  \caption{Componentes de hardware do dispositivo embarcado.}
  \label{tab:hardware}
  \begin{tabular}{p{3.5cm} p{3.5cm} p{7cm}}
    \toprule
    \textbf{Componente} & \textbf{Modelo} & \textbf{Função} \\
    \midrule
    Microcontrolador & ESP32 (Espressif) &
      Processamento central, conectividade Wi-Fi, execução do firmware, controle dos
      periféricos via GPIO. Processador dual-core, opera em 3,3 V. \\[4pt]

    Sensor de gás/fumaça & MQ-2 &
      Detecta concentrações de GLP, metano, CO e fumaça. Saída analógica lida via
      ADC no GPIO 34 do ESP32. Normalizado com 4 casas decimais. \\[4pt]

    Sensor de temp./umidade & DHT22 &
      Mede temperatura ($-40$ a $80\,^\circ$C, $\pm0,5\,^\circ$C) e umidade relativa
      (0--100\%, $\pm2\%$). Interface single-wire no GPIO 27. \\[4pt]

    Módulo relé & Relé 5V (1 canal) &
      Isola eletricamente o circuito de controle (3,3 V) do circuito de potência da
      fechadura solenoide. Acionado pelo GPIO 26 do ESP32. \\[4pt]

    Fechadura solenoide & Solenoide 12V &
      Atuador físico da barreira de acesso. Ativada pelo módulo relé mediante
      recebimento do comando \texttt{OPEN}. \\[4pt]

    LED indicador & LED 5mm (verde) &
      Sinalização visual de status da conexão e execução de comandos. GPIO 2. \\
    \bottomrule
  \end{tabular}
  \par\smallskip{\footnotesize Fonte: Próprio autor.}
\end{table}

\subsubsection{Software}

\begin{table}[H]
  \centering
  \caption{Tecnologias de software utilizadas no sistema.}
  \label{tab:software}
  \begin{tabular}{p{3cm} p{2.8cm} p{8.2cm}}
    \toprule
    \textbf{Tecnologia} & \textbf{Versão / Tipo} & \textbf{Papel no sistema} \\
    \midrule
    \textbf{C++ (Arduino)} & Arduino Framework &
      Linguagem do firmware ESP32. Bibliotecas: PubSubClient (MQTT),
      DHT (sensor DHT22), ArduinoJson (serialização JSON), WiFiClientSecure (TLS). \\[4pt]

    \textbf{Python / Django} & Django 5 + DRF &
      Backend REST API; modelos ORM; publicação MQTT via paho-mqtt;
      daemon mqtt\_subscriber para telemetria em tempo real. \\[4pt]

    \textbf{TypeScript / React} & React 19 + Vite &
      Dashboard web do anfitrião. Shadcn/UI, Tailwind CSS v4, Axios, React Router v6,
      React Hook Form + Zod. \\[4pt]

    \textbf{React Native / Expo} & Expo SDK 54 &
      Aplicativo mobile do hóspede (iOS e Android). expo-secure-store para
      persistência segura do token. \\[4pt]

    \textbf{HiveMQ Cloud} & Broker MQTT (SaaS) &
      Mensageria MQTT com TLS obrigatório na porta 8883 e autenticação
      por usuário/senha. Plano gratuito para prototipagem. \\[4pt]

    \textbf{SQLite3 / PostgreSQL} & SQLite (dev) &
      Banco de dados relacional. SQLite3 para desenvolvimento local;
      PostgreSQL para produção (migração transparente via ORM Django). \\[4pt]

    \textbf{Git / GitHub} & Monorepo &
      Controle de versão e hospedagem.
  \end{tabular}
  \par\smallskip{\footnotesize Fonte: Próprio autor.}
\end{table}

\subsection{Fatores de Risco}

\begin{table}[H]
  \centering
  \small
  \caption{Fatores de risco identificados e respectivas estratégias de mitigação.}
  \label{tab:riscos}
  \begin{tabular}{p{0.7cm} p{3.2cm} p{2.6cm} p{1.6cm} p{5.4cm}}
    \toprule
    \textbf{\#} & \textbf{Risco} & \textbf{Categoria} & \textbf{Impacto} &
    \textbf{Mitigação} \\
    \midrule

    R01 & Queda da conexão Wi-Fi no imóvel &
      Conectividade & Alto &
      Firmware trata emergências localmente sem servidor. Reconexão automática
      ao Wi-Fi em loop. \\[3pt]

    R02 & Latência elevada no broker MQTT &
      Desempenho & Médio &
      HiveMQ Cloud com SLA de disponibilidade. MQTT assíncrono; meta
      $<$3\,s monitorada. \\[3pt]

    R03 & Validação de certificado TLS ausente &
      Segurança & Médio &
      \texttt{setInsecure()} para prototipagem. Produção usará
      \texttt{setCACert()} com certificado root HiveMQ. \\[3pt]

    R04 & Clonagem de token de resgate &
      Segurança & Alto &
      Token de 12 hex (72 bits de entropia). Após resgate, status muda para
      AGENDADO, impedindo reutilização. \\[3pt]

    R05 & Imprecisão na leitura do sensor MQ-2 &
      Precisão & Médio &
      Normalização com 4 casas decimais. Limiar de alerta calibrável por
      constante no firmware. \\[3pt]

    R06 & Hardware físico indisponível para testes &
      Logístico & Alto &
      Firmware testável no simulador Wokwi. Software testável via API e
      dashboard sem hardware. \\[3pt]

    R07 & SQLite3 inadequado para produção &
      Escalabilidade & Baixo &
      ORM Django suporta migração para PostgreSQL sem alterações no código. \\[3pt]

    R08 & Status online sem WebSocket (polling) &
      Usabilidade & Baixo &
      LWT atualiza \texttt{esta\_online} via mqtt\_subscriber.
      WebSocket previsto como trabalho futuro. \\

    \bottomrule
  \end{tabular}
  \par\smallskip{\footnotesize Fonte: Próprio autor.}
\end{table}

% ============================================================
% FIM DO DOCUMENTO
% ============================================================
\end{document}
